% Autorska priprava k Realisticky.cz 020310; ulohy jsou autorske a obmenene.
% Zamerne nenahrazuje vyklad presnymi kopiemi cizich textu/prikladu.
\documentclass[11pt,a4paper]{article}
\input{../../sablony/pripravy/priprava.sty}
\usepackage{tikz}
\nastavPripravu{1. ročník}{Rovnice a nerovnice}{79}{Soustavy: dosazovací metoda}
\begin{document}
\blok{Návaznost a motivace}
V 66. hodině už studenti hledali průsečík dvou grafů. Nyní převedeme stejnou otázku do algebraického zápisu: hledáme dvojici $[x;y]$, která splní obě rovnice najednou.
\begin{definiceBox}[title={Definice - Řešení soustavy lineárních rovnic}]
\textbf{Řešení soustavy lineárních rovnic}. Je uspořádaná dvojice $[x;y]$, která současně vyhovuje všem rovnicím soustavy. Není to pouze nalezené číslo $x$.
\end{definiceBox}
\blok{Princip a odvození}
\textbf{Dosazovací metoda:} Z jedné rovnice vyjádříme jednu neznámou, vyjádřený výraz dosadíme do druhé rovnice, vyřešíme získanou rovnici, dopočítáme zbývající neznámou a prověříme obě původní rovnice.
\blok{Řešené příklady: od jednoduchého ke složitějšímu}
\begin{prikladReseniBox}[title={Příklad 1}]
Řešte $\begin{cases}x+y=8,\\x-2y=2.\end{cases}$
\tcblower
\[
\begin{aligned}
x&=8-y \\
 (8-y)-2y&=2 \\
 -3y&=-6 \\
 y&=2,\quad x=6
\end{aligned}
\]

Zkouška: $6+2=8$, $6-4=2$.
\[
\begin{aligned}
K &= \{[6;2]\}
\end{aligned}
\]
\end{prikladReseniBox}
\begin{prikladReseniBox}[title={Příklad 2}]
Řešte $\begin{cases}2x-y=4,\\x+2y=7.\end{cases}$
\tcblower
\[
\begin{aligned}
y&=2x-4 \\
 x+2(2x-4)&=7 \\
 5x&=15 \\
 x&=3,\quad y=2
\end{aligned}
\]

Zkouška $6-2=4$ a $3+4=7$.
\[
\begin{aligned}
K &= \{[3;2]\}
\end{aligned}
\]
\end{prikladReseniBox}
\begin{prikladReseniBox}[title={Příklad 3}]
Řešte $\begin{cases}7x-3y=15,\\5x+6y=27.\end{cases}$ Dosazujte ze vztahu $7x-3y=15$.
\tcblower
\[
\begin{aligned}
y&=\frac{7x-15}{3},\\
5x+6\cdot\frac{7x-15}{3}&=27,\\
5x+14x-30&=27,\\
19x&=57,\\
x&=3,\\
y&=\frac{21-15}{3}=2.
\end{aligned}
\]
Kontrola: $21-6=15$ a $15+12=27$.
\[
\begin{aligned}
K &= \{[3;2]\}
\end{aligned}
\]
\end{prikladReseniBox}

\poznamka{Žáci mají vyzkoušet sestavit jeden dosazovací postup sami. V řešení vždy piš dvojici $[x;y]$.}
\blok{Samostatné procvičování}
\begin{zadaniReseniBox}[title={Samostatná práce 1}]
$\begin{cases}x+y=9,\\x-y=1.\end{cases}$
\tcblower
\[
\begin{aligned}
x&=9-y \\
 9-2y&=1 \\
 y&=4,\quad x=5 \\[2mm]
K &= \{[5;4]\}
\end{aligned}
\]
\end{zadaniReseniBox}
\begin{zadaniReseniBox}[title={Samostatná práce 2}]
$\begin{cases}3x+y=11,\\x+y=5.\end{cases}$
\tcblower
\[
\begin{aligned}
y&=5-x \\
 3x+5-x&=11 \\
 x&=3,\quad y=2 \\[2mm]
K &= \{[3;2]\}
\end{aligned}
\]
\end{zadaniReseniBox}
\begin{zadaniReseniBox}[title={Samostatná práce 3}]
$\begin{cases}2x+y=7,\\x-y=2.\end{cases}$
\tcblower
\[
\begin{aligned}
y&=x-2 \\
 2x+x-2&=7 \\
 x&=3,\quad y=1 \\[2mm]
K &= \{[3;1]\}
\end{aligned}
\]
\end{zadaniReseniBox}
\blok{Volitelná rezerva}
\begin{zadaniReseniBox}[title={Samostatná práce 4}]
Vyřešte soustavu $\begin{cases}y=2x+1,\\x+y=10.\end{cases}$
\tcblower
\[
\begin{aligned}
x+2x+1&=10 \\
 x&=3 \\
 y&=7
\end{aligned}
\]

Zkouška: $7=2\cdot3+1$ a $3+7=10$.
\[
\begin{aligned}
K &= \{[3;7]\}
\end{aligned}
\]
\end{zadaniReseniBox}
\begin{zadaniReseniBox}[title={Rozšíření 1}]
Řešte $\begin{cases}3(x-2)+2y=x+y,\\4x+5(y+x)=3x-6.\end{cases}$
\tcblower
\[
\begin{aligned}
2x+y&=6,\\
6x+5y&=-6,\\
y&=6-2x,\\
6x+5(6-2x)&=-6,\\
-4x&=-36,\\
x&=9,\quad y=-12.
\end{aligned}
\]
Zkouška v původních rovnicích.
\[
\begin{aligned}
K &= \{[9;-12]\}
\end{aligned}
\]
\end{zadaniReseniBox}

\blok{Domácí úkol}
\begin{zadaniReseniBox}[title={Domácí úkol}]
\begin{enumerate}[label=\alph*)]
\item
$\begin{cases}x+y=10,\\x-2y=1.\end{cases}$
\item
$\begin{cases}x-y=4,\\2x+y=11.\end{cases}$
\item
$\begin{cases}3x+y=13,\\x+y=7.\end{cases}$
\item
$\begin{cases}y=2x-5,\\3x+2y=11,\end{cases}$.
\item
$\begin{cases}2x-y=3,\\3(x+y)-2x=19,\end{cases}$.
\item
$\begin{cases}3x-2y=8,\\2(x-2y)+3(2x+y)=17,\end{cases}$.
\end{enumerate}
\tcblower
\begin{enumerate}[label=\alph*)]
\item
\[
\begin{aligned}
x&=10-y \\
 10-3y&=1 \\
 y&=3,\quad x=7 \\[2mm]
K &= \{[7;3]\}
\end{aligned}
\]
\item
\[
\begin{aligned}
x&=y+4 \\
 2(y+4)+y&=11 \\
 y&=1,\quad x=5 \\[2mm]
K &= \{[5;1]\}
\end{aligned}
\]
\item
\[
\begin{aligned}
y&=7-x \\
 3x+7-x&=13 \\
 x&=3,\quad y=4 \\[2mm]
K &= \{[3;4]\}
\end{aligned}
\]
\item
\[
\begin{aligned}
3x+2(2x-5)&=11\\7x&=21\\x&=3,\quad y=1 \\[2mm]
K &= \{[3;1]\}
\end{aligned}
\]
\item
\[
\begin{aligned}
y&=2x-3\\x+3(2x-3)&=19\\7x&=28\\x&=4,\quad y=5 \\[2mm]
K &= \{[4;5]\}
\end{aligned}
\]
\item
\[
\begin{aligned}
y&=\frac{3x-8}{2}\\8x-y&=17\\16x-(3x-8)&=34\\x&=2,\quad y=-1 \\[2mm]
K &= \{[2;-1]\}
\end{aligned}
\]
\end{enumerate}
\end{zadaniReseniBox}
\poznamka{Orientační časové rozvržení: 5 + 12 + 12 + 10 + 6 = 45 min. Přednost má pochopení jádra; rezerva není součástí časového součtu. Domácí úlohy jsou určeny na samostatnou práci, řešení zůstávají učiteli.}
\end{document}
